\documentclass[conference]{IEEEtran}
\IEEEoverridecommandlockouts
% The preceding line is only needed to identify funding in the first footnote. If that is unneeded, please comment it out.
\usepackage{cite}
\usepackage{amsmath,amssymb,amsfonts}
\usepackage{algorithmic}
\usepackage{graphicx}
\usepackage{textcomp}
\usepackage{xcolor}
\def\BibTeX{{\rm B\kern-.05em{\sc i\kern-.025em b}\kern-.08em
    T\kern-.1667em\lower.7ex\hbox{E}\kern-.125emX}}
\begin{document}

\title{Conference Paper Title*\\
{\footnotesize \textsuperscript{*}Note: Sub-titles are not captured in Xplore and
should not be used}
\thanks{Identify applicable funding agency here. If none, delete this.}
}

\author{\IEEEauthorblockN{1\textsuperscript{st} Given Name Surname}
\IEEEauthorblockA{\textit{dept. name of organization (of Aff.)} \\
\textit{name of organization (of Aff.)}\\
City, Country \\
email address or ORCID}
\and
\IEEEauthorblockN{2\textsuperscript{nd} Given Name Surname}
\IEEEauthorblockA{\textit{dept. name of organization (of Aff.)} \\
\textit{name of organization (of Aff.)}\\
City, Country \\
email address or ORCID}
\and
\IEEEauthorblockN{3\textsuperscript{rd} Given Name Surname}
\IEEEauthorblockA{\textit{dept. name of organization (of Aff.)} \\
\textit{name of organization (of Aff.)}\\
City, Country \\
email address or ORCID}
\and
\IEEEauthorblockN{4\textsuperscript{th} Given Name Surname}
\IEEEauthorblockA{\textit{dept. name of organization (of Aff.)} \\
\textit{name of organization (of Aff.)}\\
City, Country \\
email address or ORCID}
\and
\IEEEauthorblockN{5\textsuperscript{th} Given Name Surname}
\IEEEauthorblockA{\textit{dept. name of organization (of Aff.)} \\
\textit{name of organization (of Aff.)}\\
City, Country \\
email address or ORCID}
\and
\IEEEauthorblockN{6\textsuperscript{th} Given Name Surname}
\IEEEauthorblockA{\textit{dept. name of organization (of Aff.)} \\
\textit{name of organization (of Aff.)}\\
City, Country \\
email address or ORCID}
}

\maketitle

\begin{abstract}
This document is a model and instructions for \LaTeX.
This and the IEEEtran.cls file define the components of your paper [title, text, heads, etc.]. *CRITICAL: Do Not Use Symbols, Special Characters, Footnotes, 
or Math in Paper Title or Abstract.
\end{abstract}

\begin{IEEEkeywords}
component, formatting, style, styling, insert
\end{IEEEkeywords}

\section{Introduction}
This document is a model and instructions for \LaTeX.
Please observe the conference page limits. 

\section{Ease of Use}

\subsection{Maintaining the Integrity of the Specifications}

The IEEEtran class file is used to format your paper and style the text. All margins, 
column widths, line spaces, and text fonts are prescribed; please do not 
alter them. You may note peculiarities. For example, the head margin
measures proportionately more than is customary. This measurement 
and others are deliberate, using specifications that anticipate your paper 
as one part of the entire proceedings, and not as an independent document. 
Please do not revise any of the current designations.

%% ------------------------- PTBR

% 2 Sensors and Technologies for Environmental Monitoring
% Sensor classification from that “sensors daily use…” - Active or passive and applications.
% 2.1 Water Monitoring
% Focus: Water level (ultrasonic, LiDAR, float sensors),Water quality (pH, turbidity, EC sensors), flood detection, river monitoring.
% Innovations: Archimedes-based sensors, object recognition methods, SAW?.

% Relevant Articles:
% 1"An Intelligent Water Level Estimation System" — Uses object recognition via deep learning for water levels. (Observations: Estimation via camera + ML.)
% 2"A Technical Evaluation of Lidar-Based Measurement of River Water Levels" — Discusses LiDAR for precise water monitoring. (High relevance.)
% 3"A low-cost ultrasonic sensor for online monitoring of liquid level" — Simple ultrasonic-based approach. (Low relevance but practical example.)
% 4"SAW Sensor based a Novel Hydrostatic Liquid Level Sensor" — Novel SAW-based liquid level sensor. (Observation: Hydrostatic pressure with SAW.)
% 5"Modeling and testing of a highly sensitive surface acoustic wave hydrostatic liquid level sensor" — Focus on SAW tech for water levels.
% 6"Development of liquid level measurement technologies: A review" — Industrial-focused review of level sensors.

% 2.2 Soil Monitoring
% Focus: Soil moisture, nutrients, agriculture precision.
% Use of SAW sensors, RFID-based systems, and nanotechnology.
% Relevant Articles:
% 1"Smart Agriculture Systems: Soil Sensors and Plant Wearables for Smart and Precision Agriculture" — Extensive review, strong base for soil monitoring section. (MIT paper, highly relevant.)
% 2"Conception and Design of WSN Sensor Nodes Based on Soil Moisture Monitoring" — Focus on soil pH, turbidity, and temperature sensors. (Observation: Detailed sensor design.)
% 3"Applications of Sensor Networks and Remote Sensing in Precision Agriculture" — General but with some relevance to WSN in agriculture. (Observation: Somewhat shallow but covers multiple angles.)

% 2.3 Air and Chemical Monitoring
% Focus: Air quality, gas detection, environmental chemical monitoring.
% "SAW Sensors for Chemical Vapors and Gases" — Discusses SAW sensor principles for air/gas detection. (Observation: Emphasizes SAW advantages like sensitivity and robustness.)
% "A Progress Review on Solid‐State LiDAR and Nanophotonic Sensors" — Discusses LiDAR, photonic sensors applicable to particulate air detection or similar use cases.

\section{Sensores e Tecnologias para Monitoramento Ambiental} \label{cap:sensors}

Nosso dia a dia é repleto de sensores, sempre coletando dados sobre o ambiente ao nosso redor. Esses sensores podem ser classificados de muitas formas diferentes, alguns exemplos de categorização relevantes para entendimento são: sensores ativos e passivos, onde os sensores ativos emitem algum tipo de sinal para medir uma grandeza física, enquanto os sensores passivos apenas captam sinais já existentes no ambiente \cite{javaid_2021_sensors}. Outra forma de classificar os sensores é a forma com que interagem com o ambiente, onde sensores de contato medem grandezas físicas diretamente em contato com o meio, enquanto sensores de não contato medem grandezas físicas sem necessidade de contato direto \cite{javaid_2021_sensors, mohindru_2023_development, wu_2023_a}. Por fim, existem ainda muitas outras formas de classificar sensores, neste trabalho optou-se por agrupar os sensores de acordo com o tipo de ambiente onde sao aplicados, como sensores de nível de água, sensores de qualidade do ar, sensores de solo, de forma a melhor guiar pesquisas nesses ambientes. 

\section{Monitoramento de Água}\label{sec:water_monitoring}

Quando falamos de monitoramento de água, podemos dividir os sensores em dois tipos principais: sensores de nível de água e sensores de qualidade da água. Os sensores de nível de água são utilizados para medir a altura da coluna d'água em corpos hídricos, como rios, lagos e reservatórios ou em ambientes controlados como tanques e poços. Já os sensores de qualidade da água são usados para medir parâmetros como pH, turbidez, condutividade elétrica (EC) e outros compostos químicos presentes na água.

\subsection{Sensores de Nível de Água}\label{subsec:water_level_sensors}

Para o monitoramento de nível de água em ambientes naturais como rios e lagos, podemos destacar os tradicionais sensores de contato, como os sensores de pressão ou linigráficos, que medem a pressão exercida pela coluna d'água sobre o sensor, um destes foi testado por \textcite{santana_2024_development}, destacando que o mesmo faz boas leituras, sendo indiferente a parametros como turbidez da água, porém o mesmo trabalho destaca o ponto fraco desta tecnologia que é a sua aplicação em ambientes hostis, visto que em situaçoes como a cheia do rio, onde há uma grande quantidade de entulhos e forte correnteza, esses sensores podem ser comprometidos. Além disso sensores de contato para monitoramento em ambientes naturais podem envolver conjuntos de boias, que detecam a altura da água em relação a um ponto fixo, ou sensores de flutuação, que utilizam um flutuador conectado a um cabo ou corda para medir a altura da água. Esses sensores são simples e eficazes, mas podem ser afetados por fatores como detritos flutuantes, variações de temperatura e corrosão \cite{mohammadrezamasoudimoghaddam_2024_a,santana_2024_development,paul_2020_a,yukawa_2025_an} .

Quando buscamos sensores para monitoramento de nível em ambientes controlados como tanques e poços, e não estamos preocupados em sensores de longo alcance, não precisamos nos preocupar tanto com fatores como detritos flutuantes, e podemos utilizar os sensores ja mencionados anteriormente para aplicações em rios e lagos, além de outros metodos mais simples como os sensores de nível de água resistivos, que medem a resistência elétrica entre dois eletrodos submersos na água. Esses sensores são econômicos e fáceis de instalar, mas podem ser afetados por corrosão e depósitos minerais \cite{santana_2024_development, mohindru_2023_development}.

Para essa categoria  monitoramento de nível em ambientes controlados existem ainda tecnologias emergentes que demonstram potencial, como os sensores passivos de nível de água baseados em ondas acústicas, o sensor SAW mede variações de deformação ou pressão na parede do tanque causadas pela mudança no nível da água, convertendo essas alterações em sinais de resposta, como os desenvolvidos por \textcite{ali_2020_saw} e \textcite{sreejith_2024_modeling}. Outra tecnologia que aparece como promissora são os sensores de nível com fibra óptica, se baseiam no principio de Arquimedes da hidrostatica e utilizam a variação da luz transmitida através de uma fibra óptica e um elemento flutuante para medir a altura da coluna d'água, o principio é demonstrado com avanços desenvolvidos por \textcite{ramos_2025_high}. Esses sensores são altamente precisos, mas ainda são complexos de instalar e operar.

Para aplicações de monitoramento sem contato, esta revisão apresenta tres destaques, sensores ultrassônicos, sensores lidar, e monitoramento remoto através de imagens geradas por aplicações aeroespaciais.

Os sensores de nível de água baseados em ultrassom, emitem ondas sonoras e medem o tempo que essas ondas levam para retornar ao sensor. Esses sensores são amplamente utilizados devido à sua precisão e capacidade de operar em ambientes com variações de temperatura e pressão \cite{mohammadrezamasoudimoghaddam_2024_a, pereira_2022_evaluation}. Por exemplo, o modelo de sensor ultrassônico GY-Us42 foi testado com resultados indicando que o erro médio do dispositivo é inferior a 3\% \cite{mohammadrezamasoudimoghaddam_2024_a}. Outro modelo, o HC-SR04, também foi avaliado como uma alternativa técnica e econômica viável para monitoramento de níveis de água \cite{pereira_2022_evaluation}, além de ser uma boa opção para educação, ciência cidadã e pesquisa devido ao seu baixo custo \cite{bresnahan_2023_a}.

Os sensores LiDAR fazem uso de ondas ópticas para medir distâncias e velocidades, sendo amplamente utilizados em metrologia, monitoramento ambiental, arqueologia e robótica\cite{behroozpour_2017_lidar, li_2022_a}. O princípio de medição do LiDAR baseia-se na rugosidade da superfície refletora para gerar reflexão não especular (ou seja, dispersão) do feixe laser emitido. A faixa de luz próxima do infravermelho (NIR) é a mais comumente utilizada para esse fim, geralmente em comprimentos de onda entre 900 e 1100 nm (270-330 THz), devido ao baixo custo dos lasers operando nessa faixa e à menor densidade de energia em comparação com o espectro visível \cite{li_2022_a, fernandezdiaz_2014_early, smart_2009_river, behroozpour_2017_lidar}. Assim como nos sensores ultrassônicos, a medição do LiDAR é feita pelo tempo de voo (TOF), mas existem dois métodos principais de medição, o TOF pulsado e o AMCW TOF, que diferem na forma como o sinal é emitido e recebido. No TOF pulsado, um pulso óptico é emitido e o tempo que leva para retornar ao sensor é medido, enquanto no AMCW TOF, uma onda contínua modulada em amplitude é utilizada, e a diferença de fase entre os sinais transmitidos e recebidos é medida para determinar a distância \cite{li_2022_a}.

Esses sensores foram explorados como uma alternativa de baixo custo para medir níveis de água a partir de pontes, com testes em laboratório e campo indicando uma boa precisão com erro de 0.1\%, mas com variações significativas causadas pela temperatura do sensor e pela rugosidade da água \cite{paul_2020_a}. Sensores LiDAR instalados em margens de rios para monitoramento de enchentes também foram testados com bons resultados, o estudo indicou que partículas suspensas na água impactam positivamente a precisão das leituras, e que o sensor poderia também ser usado para detectar a cncentração dessas partículas suspensas na água\cite{tamari_2016_flash}. Outro estudo comparou o modelo de sensor LiDAR TF-mini com sensores de pressão linigráficos, mostrando os benefícios do método de medição sem contato do LiDAR em comparação com o método de contato do linigráfico e validando o LiDAR como uma excelente escolha entre as tecnologias de medição de nível de fluido comparadas \cite{santana_2024_development}.

Outro método de monitoramento de nível de água em ambientes abertos, e a utilização de monitoramento remoto usando dados obtidos via satélite, como o trabalho realizado por \textcite{jiang_2024_monitoring}, que utiliza imagens e outros dados obtidos via satélite para estimar os níveis de uma bacia hidrográfica ou os de \textcite{ali_2024_satellite} que também usam uma tecnica similar e comparam  estimativas feitas com dados obtidos via satélite com medicoes feitas em varios pontos ao longo do rio com metodos mais tradicionais, ambos os trabalhos apresentam resultados positivos para esta tecnica.

\subsection{Qualidade da Água}
O trabalho de \textcite{ferreira_2023_conception} apresenta um projeto de monitoramento de qualidade da água, que utiliza sensores de pH, turbidez e condutividade elétrica (EC) e integra técnicas de fusão de dados locais e distribuídos com recursos de machine learning para melhorar a detecção de poluentes em tempo real. 

\section{Monitoramento do Solo}
% Focus: Soil moisture, nutrients, agriculture precision.
% Use of SAW sensors, RFID-based systems, and nanotechnology.
% Relevant Articles:
% 1"Smart Agriculture Systems: Soil Sensors and Plant Wearables for Smart and Precision Agriculture" — Extensive review, strong base for soil monitoring section. (MIT paper, highly relevant.)
% 2"Conception and Design of WSN Sensor Nodes Based on Soil Moisture Monitoring" — Focus on soil pH, turbidity, and temperature sensors. (Observation: Detailed sensor design.)
% 3"Applications of Sensor Networks and Remote Sensing in Precision Agriculture" — General but with some relevance to WSN in agriculture. (Observation: Somewhat shallow but covers multiple angles.)

O monitoramento do solo é uma ferramenta essencial para otimizar o crescimento das culturas, melhorar a eficiência produtiva e promover práticas agrícolas mais sustentáveis. Os sensores de solo permitem a medição contínua de parâmetros físicos e químicos, como umidade e concentração de nutrientes, fornecendo dados em tempo real para apoiar a tomada de decisão no campo. A demanda por essas tecnologias vem crescendo, impulsionada pelo aumento populacional, pela necessidade de ampliar a produção de alimentos e pela pressão por práticas agrícolas mais eficientes e conscientes em relação ao meio ambiente. Um exemplo representativo é o mercado de sensores de umidade do solo, que movimentou cerca de US\$ 147{,}5 milhões em 2020, com expectativa de atingir US\$ 360{,}9 milhões até 2027, refletindo o interesse global em tecnologias de agricultura digital \cite{yin_2021_smart} .

Historicamente, as recomendações de manejo agrícola foram desenvolvidas seguindo uma lógica de zonas agroecológicas amplas, como ocorreu durante a Revolução Verde, quando o foco estava apenas no aumento da produtividade por meio de fertilizantes sintéticos, sem considerar adequadamente as condições locais de solo e água, nem os impactos ambientais associados. Grande parte desse legado ainda persiste, com práticas baseadas em procedimentos centralizados e relações empíricas genéricas entre nutrientes, doses de fertilizantes e produtividade. Nesse contexto, os sensores de solo surgem como uma ferramenta fundamental para romper com esse modelo “top-down” e viabilizar uma abordagem “bottom-up”, em que as decisões de manejo passam a ser orientadas por dados reais, específicos de cada microambiente agrícola, no espaço e no tempo \cite{viscarrarossel_2016_soil}.

\subsection{SAW, RFID, e nanotechnology}
% monitoring system consists of RFID sensor, patrol car, farmland monitoring center and cloud platform
Como sabemos a agricultura pode cobrir vastas areas e nestes casos WSNs e tradiconalmente LoRaWAN são muito aplicáveis pois permitem cobrir estas áreas de forma remota \textcite{deng_2020_novel}. Dito isso, outras alternativas de cobertura para grandes areas tem sido exploradas além do mais tradicional LoRaWAN, o trabalho de \textcite{deng_2020_novel} apresenta um sistema de monitoramento que consiste em sensores RFID no solo, um veículo que se movimenta pela área monitorada coletando dados, e um centro de coleta e tratamento de dados, capaz de cobrir amplas áreas. Outro trabalho desenvolvido por ESSE CARA \textcite{akhileshnagpure_2022_water}, apresenta a utilização de drones de maneira similar a citada, para captar dados em áreas remotas.

O trabalho de \textcite{boada_2018_batteryless} apresenta o desenvolvimento de um sensor inovador de umidade do solo que opera sem bateria, utilizando tecnologia NFC com colheita de energia. O dispositivo é alimentado pelo campo magnético gerado pelo leitor NFC e realiza a medição de temperatura, umidade relativa e conteúdo volumétrico de água no solo. Um microcontrolador integrado processa os dados coletados e os transmite ao chip NFC via I2C, armazenando as informações no formato NDEF para posterior leitura. O estudo também compara diferentes métodos de medição de umidade do solo, selecionando a abordagem mais adequada às limitações energéticas do sistema. Com um princípio de funcionamento similar ao RFID, sensores passivos do tipo SAW também são apresentados como alternativas para monitoramento de características do solo, como explorado por BLABLABLA \textcite{akhileshnagpure_2022_water}. Uma forma inovadora de detecção de umidade do solo foi explorada por BLABLABLA \textcite{akhileshnagpure_2022_water} utilizando fibras ópticas instaladas em vastas áreas para verificar a concentração de água no solo.

AQUI UM PARAGRAFO SOBRE PLANTAS E NUTRIENTES, EXPLICAR SAW SENSOR E RFID?

\section{Monitoramento do Ar e Composição Química do Ar}
% Focus: Air quality, gas detection, environmental chemical monitoring.
% "SAW Sensors for Chemical Vapors and Gases" — Discusses SAW sensor principles for air/gas detection. (Observation: Emphasizes SAW advantages like sensitivity and robustness.)
% "A Progress Review on Solid‐State LiDAR and Nanophotonic Sensors" — Discusses LiDAR, photonic sensors applicable to particulate air detection or similar use cases.
Quando pensamos em ambientes urbanos integrados com IoT, dentro do conceito de smart cities, um dos parâmetros mais relevantes é a qualidade do ar, devido aos impactos diretos na saúde pública, no meio ambiente e na economia global. A poluição atmosférica em áreas urbanas, com distribuição espacial e temporal não uniforme, reforça a necessidade de sistemas de monitoramento com alta resolução espaço-temporal, algo que os sistemas tradicionais de monitoramento ainda não conseguem oferecer de forma escalável e com ampla cobertura de dados \cite{yi_2015_a}.

Nesse contexto, o avanço das tecnologias de sensores, como MEMS e redes de sensores sem fio (WSN), tem impulsionado o desenvolvimento do conceito de The Next Generation Air Pollution Monitoring System (TNGAPMS). Para esse tipo de aplicação, os gases de maior preocupação são monóxido de carbono (CO), dióxido de nitrogênio (NO$_2$), ozônio ao nível do solo (O$_3$) e dióxido de enxofre (SO$_2$). Atualmente, os sensores mais utilizados e adequados para o monitoramento de tais gases em cenários urbanos e industriais são os sensores eletroquímicos e os sensores de estado sólido (semicondutores), embora também existam outras tecnologias de baixo custo, como os sensores catalíticos, NDIR e PID, que são amplamente aplicados em diferentes contextos de detecção de gases \cite{yi_2015_a}.

\subsection{Detecção de Elementos Químicos}

O estudo de \textcite{devkota_2017_saw} investiga sensores SAW para detecção passiva de gases e vapores químicos no ar, baseados na interação dos compostos com a antena, que altera o sinal acústico recebido. Sensores SAW mostram-se viáveis para detectar gases inorgânicos como NH$_3$, NO$_2$, SO$_2$, CO, vapores orgânicos como tolueno, etanol, acetona e agentes químicos de guerra.
O trabalho destaca a importância da estabilidade dos materiais em ambientes extremos, o uso de antenas para operação sem fio e baixo consumo energético, e aponta para futuras pesquisas em materiais sensíveis, sensores flexíveis e arranjos multi-elemento para monitoramento simultâneo de vários gases.

FALAR DE OUTROS SENSORES MAIS COMUNS PARA ESTA APL

\subsection{Poluição e Qualidade do Ar}
O trabalho de \textcite{karagulian_2019_review} analisa o desempenho de sensores de baixo custo (LCS) para o monitoramento de diversos poluentes atmosféricos, incluindo monóxido de carbono (CO), óxidos de nitrogênio (NO e NO$_2$), ozônio (O$_3$) e material particulado (PM$_{2.5}$). A revisão destaca o potencial desses sensores para ampliar a cobertura espacial em áreas urbanas e remotas, além de avaliar diferentes métodos de calibração, como MLR, ANN, SVR e RF, levando em conta fatores como a umidade relativa, que influencia significativamente a medição de partículas.

O trabalho \textcite{yi_2015_a} apresenta uma revisão dos sistemas de monitoramento da poluição do ar baseados em redes de sensores sem fio (WSN), classificando-os em três categorias principais: Redes de Sensores Estáticas (SSN), Redes de Sensores Comunitárias (CSN) e Redes de Sensores Veiculares (VSN), de acordo com os tipos de portadores dos sensores. A análise destaca que muitas soluções atuais já demonstram ser viáveis em termos de resolução espaço-temporal, custo, eficiência energética, facilidade de implantação, manutenção e acesso público aos dados. No entanto, permanecem desafios como a falta de aquisição de dados tridimensionais, limitações na capacidade de monitoramento ativo e o uso de portadores não controlados ou parcialmente controlados, aspectos que devem ser aprimorados nas futuras gerações de sistemas de monitoramento da poluição do ar.

%% ------------------------- PTBR
\section{Prepare Your Paper Before Styling}
Before you begin to format your paper, first write and save the content as a 
separate text file. Complete all content and organizational editing before 
formatting. Please note sections \ref{AA}--\ref{SCM} below for more information on 
proofreading, spelling and grammar.

Keep your text and graphic files separate until after the text has been 
formatted and styled. Do not number text heads---{\LaTeX} will do that 
for you.

\subsection{Abbreviations and Acronyms}\label{AA}
Define abbreviations and acronyms the first time they are used in the text, 
even after they have been defined in the abstract. Abbreviations such as 
IEEE, SI, MKS, CGS, ac, dc, and rms do not have to be defined. Do not use 
abbreviations in the title or heads unless they are unavoidable.

\subsection{Units}
\begin{itemize}
\item Use either SI (MKS) or CGS as primary units. (SI units are encouraged.) English units may be used as secondary units (in parentheses). An exception would be the use of English units as identifiers in trade, such as ``3.5-inch disk drive''.
\item Avoid combining SI and CGS units, such as current in amperes and magnetic field in oersteds. This often leads to confusion because equations do not balance dimensionally. If you must use mixed units, clearly state the units for each quantity that you use in an equation.
\item Do not mix complete spellings and abbreviations of units: ``Wb/m\textsuperscript{2}'' or ``webers per square meter'', not ``webers/m\textsuperscript{2}''. Spell out units when they appear in text: ``. . . a few henries'', not ``. . . a few H''.
\item Use a zero before decimal points: ``0.25'', not ``.25''. Use ``cm\textsuperscript{3}'', not ``cc''.)
\end{itemize}

\subsection{Equations}
Number equations consecutively. To make your 
equations more compact, you may use the solidus (~/~), the exp function, or 
appropriate exponents. Italicize Roman symbols for quantities and variables, 
but not Greek symbols. Use a long dash rather than a hyphen for a minus 
sign. Punctuate equations with commas or periods when they are part of a 
sentence, as in:
\begin{equation}
a+b=\gamma\label{eq}
\end{equation}

Be sure that the 
symbols in your equation have been defined before or immediately following 
the equation. Use ``\eqref{eq}'', not ``Eq.~\eqref{eq}'' or ``equation \eqref{eq}'', except at 
the beginning of a sentence: ``Equation \eqref{eq} is . . .''

\subsection{\LaTeX-Specific Advice}

Please use ``soft'' (e.g., \verb|\eqref{Eq}|) cross references instead
of ``hard'' references (e.g., \verb|(1)|). That will make it possible
to combine sections, add equations, or change the order of figures or
citations without having to go through the file line by line.

Please don't use the \verb|{eqnarray}| equation environment. Use
\verb|{align}| or \verb|{IEEEeqnarray}| instead. The \verb|{eqnarray}|
environment leaves unsightly spaces around relation symbols.

Please note that the \verb|{subequations}| environment in {\LaTeX}
will increment the main equation counter even when there are no
equation numbers displayed. If you forget that, you might write an
article in which the equation numbers skip from (17) to (20), causing
the copy editors to wonder if you've discovered a new method of
counting.

{\BibTeX} does not work by magic. It doesn't get the bibliographic
data from thin air but from .bib files. If you use {\BibTeX} to produce a
bibliography you must send the .bib files. 

{\LaTeX} can't read your mind. If you assign the same label to a
subsubsection and a table, you might find that Table I has been cross
referenced as Table IV-B3. 

{\LaTeX} does not have precognitive abilities. If you put a
\verb|\label| command before the command that updates the counter it's
supposed to be using, the label will pick up the last counter to be
cross referenced instead. In particular, a \verb|\label| command
should not go before the caption of a figure or a table.

Do not use \verb|\nonumber| inside the \verb|{array}| environment. It
will not stop equation numbers inside \verb|{array}| (there won't be
any anyway) and it might stop a wanted equation number in the
surrounding equation.

\subsection{Some Common Mistakes}\label{SCM}
\begin{itemize}
\item The word ``data'' is plural, not singular.
\item The subscript for the permeability of vacuum $\mu_{0}$, and other common scientific constants, is zero with subscript formatting, not a lowercase letter ``o''.
\item In American English, commas, semicolons, periods, question and exclamation marks are located within quotation marks only when a complete thought or name is cited, such as a title or full quotation. When quotation marks are used, instead of a bold or italic typeface, to highlight a word or phrase, punctuation should appear outside of the quotation marks. A parenthetical phrase or statement at the end of a sentence is punctuated outside of the closing parenthesis (like this). (A parenthetical sentence is punctuated within the parentheses.)
\item A graph within a graph is an ``inset'', not an ``insert''. The word alternatively is preferred to the word ``alternately'' (unless you really mean something that alternates).
\item Do not use the word ``essentially'' to mean ``approximately'' or ``effectively''.
\item In your paper title, if the words ``that uses'' can accurately replace the word ``using'', capitalize the ``u''; if not, keep using lower-cased.
\item Be aware of the different meanings of the homophones ``affect'' and ``effect'', ``complement'' and ``compliment'', ``discreet'' and ``discrete'', ``principal'' and ``principle''.
\item Do not confuse ``imply'' and ``infer''.
\item The prefix ``non'' is not a word; it should be joined to the word it modifies, usually without a hyphen.
\item There is no period after the ``et'' in the Latin abbreviation ``et al.''.
\item The abbreviation ``i.e.'' means ``that is'', and the abbreviation ``e.g.'' means ``for example''.
\end{itemize}
An excellent style manual for science writers is \cite{b7}.

\subsection{Authors and Affiliations}
\textbf{The class file is designed for, but not limited to, six authors.} A 
minimum of one author is required for all conference articles. Author names 
should be listed starting from left to right and then moving down to the 
next line. This is the author sequence that will be used in future citations 
and by indexing services. Names should not be listed in columns nor group by 
affiliation. Please keep your affiliations as succinct as possible (for 
example, do not differentiate among departments of the same organization).

\subsection{Identify the Headings}
Headings, or heads, are organizational devices that guide the reader through 
your paper. There are two types: component heads and text heads.

Component heads identify the different components of your paper and are not 
topically subordinate to each other. Examples include Acknowledgments and 
References and, for these, the correct style to use is ``Heading 5''. Use 
``figure caption'' for your Figure captions, and ``table head'' for your 
table title. Run-in heads, such as ``Abstract'', will require you to apply a 
style (in this case, italic) in addition to the style provided by the drop 
down menu to differentiate the head from the text.

Text heads organize the topics on a relational, hierarchical basis. For 
example, the paper title is the primary text head because all subsequent 
material relates and elaborates on this one topic. If there are two or more 
sub-topics, the next level head (uppercase Roman numerals) should be used 
and, conversely, if there are not at least two sub-topics, then no subheads 
should be introduced.

\subsection{Figures and Tables}
\paragraph{Positioning Figures and Tables} Place figures and tables at the top and 
bottom of columns. Avoid placing them in the middle of columns. Large 
figures and tables may span across both columns. Figure captions should be 
below the figures; table heads should appear above the tables. Insert 
figures and tables after they are cited in the text. Use the abbreviation 
``Fig.~\ref{fig}'', even at the beginning of a sentence.

\begin{table}[htbp]
\caption{Table Type Styles}
\begin{center}
\begin{tabular}{|c|c|c|c|}
\hline
\textbf{Table}&\multicolumn{3}{|c|}{\textbf{Table Column Head}} \\
\cline{2-4} 
\textbf{Head} & \textbf{\textit{Table column subhead}}& \textbf{\textit{Subhead}}& \textbf{\textit{Subhead}} \\
\hline
copy& More table copy$^{\mathrm{a}}$& &  \\
\hline
\multicolumn{4}{l}{$^{\mathrm{a}}$Sample of a Table footnote.}
\end{tabular}
\label{tab1}
\end{center}
\end{table}

\begin{figure}[htbp]
\centerline{\includegraphics{fig1.png}}
\caption{Example of a figure caption.}
\label{fig}
\end{figure}

Figure Labels: Use 8 point Times New Roman for Figure labels. Use words 
rather than symbols or abbreviations when writing Figure axis labels to 
avoid confusing the reader. As an example, write the quantity 
``Magnetization'', or ``Magnetization, M'', not just ``M''. If including 
units in the label, present them within parentheses. Do not label axes only 
with units. In the example, write ``Magnetization (A/m)'' or ``Magnetization 
\{A[m(1)]\}'', not just ``A/m''. Do not label axes with a ratio of 
quantities and units. For example, write ``Temperature (K)'', not 
``Temperature/K''.

\section*{Acknowledgment}

The preferred spelling of the word ``acknowledgment'' in America is without 
an ``e'' after the ``g''. Avoid the stilted expression ``one of us (R. B. 
G.) thanks $\ldots$''. Instead, try ``R. B. G. thanks$\ldots$''. Put sponsor 
acknowledgments in the unnumbered footnote on the first page.

\section*{References}

Please number citations consecutively within brackets \cite{b1}. The 
sentence punctuation follows the bracket \cite{b2}. Refer simply to the reference 
number, as in \cite{b3}---do not use ``Ref. \cite{b3}'' or ``reference \cite{b3}'' except at 
the beginning of a sentence: ``Reference \cite{b3} was the first $\ldots$''

Number footnotes separately in superscripts. Place the actual footnote at 
the bottom of the column in which it was cited. Do not put footnotes in the 
abstract or reference list. Use letters for table footnotes.

Unless there are six authors or more give all authors' names; do not use 
``et al.''. Papers that have not been published, even if they have been 
submitted for publication, should be cited as ``unpublished'' \cite{b4}. Papers 
that have been accepted for publication should be cited as ``in press'' \cite{b5}. 
Capitalize only the first word in a paper title, except for proper nouns and 
element symbols.

For papers published in translation journals, please give the English 
citation first, followed by the original foreign-language citation \cite{b6}.

\begin{thebibliography}{00}
\bibitem{b1} G. Eason, B. Noble, and I. N. Sneddon, ``On certain integrals of Lipschitz-Hankel type involving products of Bessel functions,'' Phil. Trans. Roy. Soc. London, vol. A247, pp. 529--551, April 1955.
\bibitem{b2} J. Clerk Maxwell, A Treatise on Electricity and Magnetism, 3rd ed., vol. 2. Oxford: Clarendon, 1892, pp.68--73.
\bibitem{b3} I. S. Jacobs and C. P. Bean, ``Fine particles, thin films and exchange anisotropy,'' in Magnetism, vol. III, G. T. Rado and H. Suhl, Eds. New York: Academic, 1963, pp. 271--350.
\bibitem{b4} K. Elissa, ``Title of paper if known,'' unpublished.
\bibitem{b5} R. Nicole, ``Title of paper with only first word capitalized,'' J. Name Stand. Abbrev., in press.
\bibitem{b6} Y. Yorozu, M. Hirano, K. Oka, and Y. Tagawa, ``Electron spectroscopy studies on magneto-optical media and plastic substrate interface,'' IEEE Transl. J. Magn. Japan, vol. 2, pp. 740--741, August 1987 [Digests 9th Annual Conf. Magnetics Japan, p. 301, 1982].
\bibitem{b7} M. Young, The Technical Writer's Handbook. Mill Valley, CA: University Science, 1989.
\end{thebibliography}
\vspace{12pt}
\color{red}
IEEE conference templates contain guidance text for composing and formatting conference papers. Please ensure that all template text is removed from your conference paper prior to submission to the conference. Failure to remove the template text from your paper may result in your paper not being published.

\end{document}
